\section{Conclusion}
\label{sec:conclusion}

We tested four quantitative claims about system-prompt compliance with a controlled benchmark of 40 verifiable rules and 4 conflict pairs on 9 open models.
Compliance degrades continuously.
It falls by about 7.7\pp (15\% relative) per doubling of prompt length at a fixed rule count, declines with every rule added from about five rules on, and drops by 30--83\% when rules contradict.
All of these effects are steeper than the hypothesized magnitudes, and the predicted size thresholds of 10 and 20 rules did not appear.
Model recipe mattered more than parameter count.

The key takeaway is about placement: in long system prompts, put the rules that must be followed in one block at the end; in short rule lists, put the most important rules first.

Several directions follow.
We plan to repeat the study with frontier API models and with bf16 versions of the 3B, 8B, and 14B models; to use realistic rule types such as tone, refusal, and tool-use policies, scored by an LLM judge; to fully cross placement, count, and length, including repeating critical rules at the end; to study multi-turn sessions; and to test whether explicit priority markers or a stated conflict-resolution order prevent compromises that satisfy neither rule.
